\documentclass{beamer}
\usepackage{ctex, hyperref}
\usepackage[T1]{fontenc}

% other packages
\usepackage{latexsym,amsmath,xcolor,multicol,booktabs,calligra}
\usepackage{graphicx,pstricks,listings,stackengine}
\usepackage{UCAS}
\usepackage{tikz}
\usetikzlibrary{positioning, arrows.meta, shapes.geometric}

% defs
\def\cmd#1{\texttt{\color{red}\footnotesize $\backslash$#1}}
\def\env#1{\texttt{\color{blue}\footnotesize #1}}
\definecolor{deepblue}{rgb}{0,0,0.5}
\definecolor{deepred}{rgb}{0.6,0,0}
\definecolor{deepgreen}{rgb}{0,0.5,0}
\definecolor{halfgray}{gray}{0.55}
\definecolor{warnorange}{RGB}{230,126,34}

\lstset{
    basicstyle=\ttfamily\small,
    keywordstyle=\bfseries\color{deepblue},
    emphstyle=\ttfamily\color{deepred},
    stringstyle=\color{deepgreen},
    numbers=left,
    numberstyle=\small\color{halfgray},
    rulesepcolor=\color{red!20!green!20!blue!20},
    frame=shadowbox,
}

\author[J. Luo]{Junyuan Luo \and Maverick Team}
\title[Maverick-SORT: Final Pitch]{Smart Warehouse Assortment Optimization}
\subtitle{Answering the Questions That Matter --- Pitch \& Demo Day}
\institute[SDC]{Sino-Danish College, UCAS}
\date{10th June, 2026}

% Global font size reduction for content
\setbeamerfont{frametitle}{size=\normalsize}
\setbeamerfont{framesubtitle}{size=\small}

\begin{document}

\kaishu
\begin{frame}
    \titlepage
    \begin{figure}[htpb]
        \begin{center}
            \includegraphics[width=0.15\linewidth]{pic/CAS_Logo.png}
        \end{center}
    \end{figure}
\end{frame}

%------------------------------------------------
\begin{frame}{Outline}
    \small
    \tableofcontents[sectionstyle=show,subsectionstyle=show/shaded/hide,subsubsectionstyle=show/shaded/hide]
\end{frame}

%=================================================
\section{1. Product \& User}
%=================================================

\begin{frame}{What Problem, and for Whom?}
  \framesubtitle{Be Specific About the Person, Not the Market}
  \small

  \textbf{The Person:} Li Wei, Warehouse Operations Manager at a tier-1 pharmaceutical distributor in Jiangsu Province. He oversees 90,000+ orders/day across 128 logistics centres, with 340 pickers (60\% temporary during peak seasons).

  \vspace{0.25cm}
  \textbf{The Problem:} Every morning at 06:30, Li faces a single decision:
  \textit{``Which orders go into which wave, in what sequence?''}

  \vspace{0.2cm}
  \begin{columns}
    \begin{column}{0.48\textwidth}
      \textbf{Current Reality:}
      \begin{itemize}
        \item Rules-of-thumb on whiteboards
        \item 38\% of SKUs are temperature-sensitive
        \item 0.35\% error rate = RMB 9M/year in returns
      \end{itemize}
    \end{column}
    \begin{column}{0.48\textwidth}
      \textbf{Cost of Getting It Wrong:}
      \begin{itemize}
        \item Cold-chain breach: batch recall, GSP fine
        \item Missed deadline: hospital penalty
        \item Over-staffing peaks: 2.5x labour cost surge
      \end{itemize}
    \end{column}
  \end{columns}
\end{frame}

\begin{frame}{Feature 1: KGDRL Core Scheduler}
  \framesubtitle{Why We Built This First}
  \small

  We deliberately built only two things to production depth --- everything else is scaffolding.

  \vspace{0.3cm}
  \textbf{\color{UBCblue}Feature 1: KGDRL Core Scheduler}
  \begin{itemize}
    \item \textbf{What it does:} Dynamically groups orders into waves using Knowledge-Graph-Guided Deep RL (PPO + GAT + KL constraint).
    \item \textbf{Why this feature:} Without the algorithm, we have nothing to sell. Every other page in the dashboard is decoration if the routing logic is wrong.
  \end{itemize}

  \vspace{0.2cm}
  \textbf{Key insight:} The KG encodes pharma-domain rules (e.g., ``frozen products must ship before 10 AM'') directly into the policy via KL-divergence constraints.
\end{frame}

\begin{frame}{Feature 2: Streamlit Pro Dashboard}
  \framesubtitle{Built for Trust, Not for Demo}
  \small

  \textbf{\color{UBCblue}Feature 2: Streamlit Pro Decision Dashboard}
  \begin{itemize}
    \item \textbf{What it does:} Real-time Command Center + What-If Simulator + Strategy Optimizer + ROI Calculator.
    \item \textbf{Why this feature:} Li Wei will not trust a black-box API. He needs to \textit{see} the wave form, test scenarios, and calculate payback before asking his CFO for budget.
  \end{itemize}

  \vspace{0.2cm}
  \textbf{Product positioning:}
  \begin{center}
    \resizebox{0.75\textwidth}{!}{
    \begin{tikzpicture}[
      node distance=0.5cm,
      box/.style={rectangle, rounded corners, draw=UBCblue, fill=UBCblue!8, thick, text width=3.2cm, align=center, minimum height=0.7cm, font=\small}
    ]
      \node[box] (a) {Do the work? \color{red}\textbf{No}};
      \node[box, below=of a] (b) {Help do it better? \color{warnorange}\textbf{Partially}};
      \node[box, below=of b] (c) {Reveal the invisible? \color{deepgreen}\textbf{Yes}};
    \end{tikzpicture}
    }
  \end{center}

  \textit{Validation:} Li Wei spent 70\% of interview time complaining about ``not knowing if today's schedule was optimal.'' He spent 0\% asking for full automation.
\end{frame}

%=================================================
\section{2. Value \& Stakeholders}
%=================================================

\begin{frame}{Who Uses, Who Pays, Who Decides?}
  \framesubtitle{Three Different People, and the Tension Between Them}
  \small

  \begin{center}
    \resizebox{0.85\textwidth}{!}{
    \begin{tikzpicture}[
      node distance=1cm and 2cm,
      person/.style={rectangle, rounded corners, draw=UBCblue, fill=UBCblue!5, thick, text width=2.8cm, align=center, minimum height=1.4cm, font=\small},
      arrow/.style={thick, ->, >=Stealth, draw=halfgray}
    ]
      \node[person] (user) {\textbf{User}\\Li Wei\\(Ops Manager)};
      \node[person, right=of user] (payer) {\textbf{Payer}\\CFO / Procurement\\(RMB 8,999/mo)};
      \node[person, right=of payer] (decider) {\textbf{Decider}\\VP Operations\\(Adoption sign-off)};

      \draw[arrow] (user) -- node[above, font=\footnotesize] {``Saves me headaches''} (payer);
      \draw[arrow] (payer) -- node[above, font=\footnotesize] {``Show me ROI''} (decider);
      \draw[arrow, bend left=20] (user) to node[below, font=\footnotesize] {``I need it''} (decider);
    \end{tikzpicture}
    }
  \end{center}

  \vspace{0.2cm}
  \textbf{When they disagree, whose view wins?}
  \begin{itemize}
    \item The CFO wins on pricing. Two tiers: Essential (RMB 2,999/mo) for trial, Pro (RMB 8,999/mo) for scale.
    \item \textbf{What this costs the product:} We cannot give Li Wei everything he wants (e.g., full ERP integration) because the CFO will not pay until ROI is proven.
  \end{itemize}
\end{frame}

\begin{frame}{What Changes in Someone's Day?}
  \framesubtitle{And Is That Change Worth the Cost?}
  \small

  \begin{columns}
    \begin{column}{0.48\textwidth}
      \textbf{Before Maverick:}
      \begin{itemize}
        \item 06:30--07:30: Manual wave planning (1 hour)
        \item 14:00: Peak surge; calls agency for 50 temp workers (reactive)
        \item 17:00: No data on whether schedule was optimal
      \end{itemize}
    \end{column}
    \begin{column}{0.48\textwidth}
      \textbf{With Maverick:}
      \begin{itemize}
        \item 06:30--06:35: Review AI plan, adjust if needed (5 min)
        \item 14:00: System shrank capacity 30 min ago (proactive)
        \item 17:00: Dashboard shows cost, SLA, compliance vs. benchmark
      \end{itemize}
    \end{column}
  \end{columns}

  \vspace{0.25cm}
  \textbf{Worth the cost?} At RMB 8,999/mo per warehouse, break-even is \textbf{month 5} under neutral assumptions (15\% efficiency gain). At RMB 9M/year in picking-error losses, a 10\% reduction pays for the software twice over.
\end{frame}

\begin{frame}{What Is Hard to Copy?}
  \framesubtitle{Where Does Our Own Judgment Sit?}
  \small

  \textbf{Anyone can build a rough version:} PPO + Streamlit + simulated data. The tools are open-source; the architecture is textbook.

  \vspace{0.25cm}
  \textbf{What is hard to copy:}
  \begin{enumerate}
    \item \textbf{Pharma-specific knowledge graph:} GSP temperature-segregation rules, FEFO priority logic, bimodal arrival patterns. A general AI company does not have this.
    \item \textbf{Human-in-the-loop trust design:} The What-If Simulator exists because we judged that Chinese warehouse managers will not accept AI recommendations they cannot challenge. That is a \textit{cultural} insight, not an algorithmic one.
    \item \textbf{The ablation protocol:} We spent 3 days validating KGDRL vs. vanilla PPO under pharma constraints. That judgment --- what to measure, what ``good enough'' means --- is not in the code.
  \end{enumerate}
\end{frame}

%=================================================
\section{3. Adoption}
%=================================================

\begin{frame}{The Binding Constraint}
  \framesubtitle{TOE Framework Applied Honestly}
  \small

  \begin{center}
    \resizebox{0.8\textwidth}{!}{
    \begin{tikzpicture}[
      node distance=0.8cm and 1.2cm,
      box/.style={rectangle, rounded corners, thick, text width=2.8cm, align=center, minimum height=1cm, font=\small},
      tech/.style={box, draw=deepgreen, fill=deepgreen!8},
      org/.style={box, draw=deepred, fill=deepred!8},
      env/.style={box, draw=UBCblue, fill=UBCblue!8}
    ]
      \node[tech] (t) {\textbf{Technology}\\\color{deepgreen}\textbf{Resolved}};
      \node[org, right=of t] (o) {\textbf{Organization}\\\color{deepred}\textbf{Binding}};
      \node[env, right=of o] (e) {\textbf{Environment}\\\color{UBCblue}\textbf{Catalyst}};

      \node[below=0.3cm of t, font=\footnotesize, text width=2.8cm, align=center] {KGDRL + Streamlit runs end-to-end};
      \node[below=0.3cm of o, font=\footnotesize, text width=2.8cm, align=center] {Temp worker turnover; TAM learning curve};
      \node[below=0.3cm of e, font=\footnotesize, text width=2.8cm, align=center] {GSP fines force digital adoption};
    \end{tikzpicture}
    }
  \end{center}

  \vspace{0.2cm}
  \textbf{Would resolving this unlock adoption, or reveal the next constraint?}

  \textit{Honest answer:} Resolving worker resistance would reveal the \textbf{data integration} constraint. Li Wei's WMS is a legacy system from 2014. Even if he trusts the algorithm, someone must build the API bridge. That is why we priced the Pro tier to include API-first integration --- it is not a feature, it is an admission that adoption requires plumbing, not just AI.
\end{frame}

\begin{frame}{First-Minute Trust \& Realistic Early Adopter}
  \framesubtitle{Is It Useful Enough, and Easy Enough?}
  \small

  \textbf{First-minute trust test:}
  \begin{itemize}
    \item Li Wei opens the dashboard. He sees a 3D Profit Mountain he can rotate.
    \item He drags the Wave Capacity slider from 20 to 15. Cost updates in real time.
    \item \textbf{Within 60 seconds, he understands the system responds to his input.}
  \end{itemize}

  \vspace{0.2cm}
  \textbf{Realistic early adopter:}
  \begin{itemize}
    \item \textbf{Who:} Regional pharma distributors with 1--3 warehouses, RMB 50--200M revenue, currently on Excel + whiteboard.
    \item \textbf{Why they try it:} Small enough for CFO sign-off in one meeting; large enough that 90,000 orders/day is aspirational but believable.
    \item \textbf{What must change before mainstream follows:} \textit{A named customer case study.} Until then, SAP EWM has the reference advantage.
  \end{itemize}
\end{frame}

%=================================================
\section{4. Building with AI}
%=================================================

\begin{frame}{Handed to AI (Claude Code)}
  \framesubtitle{What the AI Did}
  \small

  \textbf{\color{deepgreen}Handed to AI:}
  \begin{itemize}
    \item PPO network architecture (1,100 lines)
    \item Streamlit dashboard CSS and layout
    \item Multi-objective NSGA-II scaffold
    \item Translation (ZH $\leftrightarrow$ EN)
    \item Bug fixes and syntax debugging
  \end{itemize}

  \vspace{0.2cm}
  \textit{Time saved: $\sim$13 days $\rightarrow$ 4 days (3.25x)}

  \vspace{0.2cm}
  \textbf{Pattern:} The AI excelled at implementation and multilingual translation. It generated hundreds of lines of structured code in minutes.
\end{frame}

\begin{frame}{Human Judgment Did the Real Work}
  \framesubtitle{What We Decided}
  \small

  \textbf{\color{deepred}Human judgment:}
  \begin{itemize}
    \item Which pain point to solve (wave allocation vs. inventory vs. demand forecasting)
    \item Reward function weights: $\alpha_{\text{temp}} = 100$ vs. $\alpha_{\text{dist}} = 1$
    \item What ``good enough'' means: ablation protocol design
    \item Business model: per-warehouse vs. per-order pricing
    \item When to stop training (convergence criteria)
  \end{itemize}

  \vspace{0.3cm}
  \centering
  \textit{The AI wrote the code. The human decided what code was worth writing.}
\end{frame}

\begin{frame}{What We Changed Around the GenAI Model}
  \framesubtitle{Data, Context, Instructions}
  \small

  \textbf{We did not build the LLM. We built the context around it:}

  \vspace{0.15cm}
  \begin{enumerate}
    \item \textbf{Data:} Calibrated simulated orders against the casebook (55\% ambient, 25\% cool, 15\% cold, 5\% frozen; bimodal arrival at 09:00 and 14:00; peak multiplier 2.8x).
    \item \textbf{Context:} Pharma-specific knowledge graph with 4 node types and 4 edge types. Not generic RL --- RL constrained by GSP rules.
    \item \textbf{Instructions:} Prompts that translate DRL actions into natural language (``The system closed the wave at minute 3 because...'').
  \end{enumerate}

  \vspace{0.2cm}
  \textbf{How we knew it was good enough:}
  \begin{itemize}
    \item Technical: Ablation across 20 random seeds
    \item Business: Li Wei would understand the dashboard without training
    \item Academic: Reward curves converge; GAT weights correlate with zone proximity
  \end{itemize}
\end{frame}

\begin{frame}{Where AI Is Confidently Wrong}
  \framesubtitle{And What Stops a User from Trusting It There}
  \small

  \textbf{Highest-risk failure mode:} \textbf{Overfitting to simulated demand patterns.}

  \vspace{0.15cm}
  Our training assumes bimodal arrival (09:00 and 14:00 peaks). If Li Wei's warehouse sees a \textbf{unimodal} surge at 11:00 due to a hospital contract, the EWMA forecaster will under-predict for 2--3 hours before adapting.

  \vspace{0.2cm}
  \textbf{Three safeguards:}
  \begin{enumerate}
    \item \textbf{Confidence bands:} Forecast chart shows $\pm$1 SD. Li Wei sees when the model is uncertain.
    \item \textbf{Human override:} Every recommendation has a ``Reject \& Manual'' button. The system does not act autonomously.
    \item \textbf{Alert escalation:} If observed rate exceeds predicted by $>$50\% for 15 min, the Alert Center triggers a ``Model Drift'' warning.
  \end{enumerate}

  \vspace{0.15cm}
  \textit{We would not let it act without human checking --- not because we distrust AI, but because GSP compliance liability rests with the warehouse manager, not the vendor.}
\end{frame}

%=================================================
\section{5. Real Results}
%=================================================

\begin{frame}{Ablation Study: What the Numbers Actually Say}
  \framesubtitle{20 Random Seeds, 6 Methods, Honest Metrics}
  \small

  \begin{center}
    \footnotesize
    \begin{tabular}{lccccc}
      \toprule
      \textbf{Method} & \textbf{Reward} & \textbf{Waves} & \textbf{Dist.} & \textbf{Misses} & \textbf{Viol.} \\
      \midrule
      FCFS & 1,787 & 85.9 & 691 & 0.0 & 38.5 \\
      ZONE\_NN & 1,734 & 82.9 & 669 & 0.0 & 37.3 \\
      EDD & 1,752 & 83.1 & 668 & 0.0 & 37.5 \\
      TZU & 145 & 83.0 & 667 & 0.0 & 21.4 \\
      \midrule
      \textbf{Vanilla PPO} & \textbf{4,007} & 87.3 & 695 & 4.4 & 36.7 \\
      KGDRL (no KL) & 1,841 & 84.4 & 678 & \textbf{0.0} & 38.7 \\
      KGDRL (full) & 1,909 & 82.8 & 666 & \textbf{0.0} & 39.0 \\
      \bottomrule
    \end{tabular}
  \end{center}

  \vspace{0.2cm}
  \textbf{Honest reading:}
  \begin{itemize}
    \item Vanilla PPO achieves the highest reward --- but pays with \textbf{4.4 deadline misses per episode}. In pharma, one miss can mean a drug shortage.
    \item KGDRL sacrifices raw reward for \textbf{zero deadline misses} and lower variance. That is the trade-off Li Wei cares about.
    \item The gap between KGDRL and TZU is modest. Our moat is the \textbf{system} (algorithm + dashboard + trust layer).
  \end{itemize}
\end{frame}

%=================================================
\section{6. Honest Reflection}
%=================================================

\begin{frame}{What We Are Least Sure Of}
  \framesubtitle{And What We Would Do with One More Week}
  \small

  \textbf{Assumption we are least sure of:} \textit{``Warehouse managers want to \textbf{see} the algorithm's reasoning more than they want it to be perfect.''}

  \vspace{0.15cm}
  We invested heavily in the What-If Simulator and AI Copilot because of this belief. But we have not tested it with a real Li Wei. It is possible that what he actually wants is simply a lower price and a SAP integration certificate.

  \vspace{0.25cm}
  \textbf{With one more week:}
  \begin{enumerate}
    \item \textbf{Day 1--2:} Run 3 structured interviews with actual pharma warehouse managers (not the casebook persona).
    \item \textbf{Day 3--4:} Build the FastAPI backend bridge to ingest real WMS data formats.
    \item \textbf{Day 5--7:} Conduct a live A/B test: one shift with human scheduling, one with KGDRL recommendations, measuring picking distance and error rate.
  \end{enumerate}
\end{frame}

\begin{frame}{What Building This Taught Us}
  \framesubtitle{About AI, Innovation, and What We Still Do Not Understand}
  \small

  \textbf{What we did not know four weeks ago:}
  \begin{enumerate}
    \item \textbf{AI accelerates implementation, not validation.} Claude generated 2,900 lines of code in one day. But designing the ablation protocol --- deciding what ``better'' means --- took three days of human debate.
    \item \textbf{The UI is the moat.} In B2B software, the algorithm that wins is not the one with the highest reward --- it is the one the operations manager trusts enough to defend to the CFO.
    \item \textbf{Simulated data is a trap.} It lets you demo beautifully, but it hides the integration cost. Every number in our ablation is ``real'' within the simulation. None of it is real in a warehouse.
  \end{enumerate}

  \vspace{0.2cm}
  \textbf{What we still do not understand:}
  \begin{itemize}
    \item Why the GAT attention weights correlate with zone proximity \textit{most of the time}, but occasionally focus on seemingly irrelevant nodes. Is the model learning something we did not encode, or is it noise?
  \end{itemize}
\end{frame}

%=================================================
\section{7. Closing}
%=================================================

\begin{frame}{Summary: What Maverick Is, and Is Not}
  \framesubtitle{A Realistic Closing for Demo Day}
  \small

  \begin{columns}
    \begin{column}{0.48\textwidth}
      \textbf{Maverick \color{deepgreen}is:}
      \begin{itemize}
        \item A \textbf{decision-support system}, not a replacement for human judgment.
        \item Built on \textbf{KGDRL}, with honest ablation data showing trade-offs.
        \item Packaged in a \textbf{dashboard designed for trust}.
        \item Priced to \textbf{pay for itself in month 5} under conservative assumptions.
      \end{itemize}
    \end{column}
    \begin{column}{0.48\textwidth}
      \textbf{Maverick \color{deepred}is not:}
      \begin{itemize}
        \item A \textbf{patented product} (substantive examination pending).
        \item \textbf{Trained on real warehouse data} (all results are simulation-based).
        \item \textbf{Autonomous} (human check is mandatory).
      \end{itemize}
    \end{column}
  \end{columns}

  \vspace{0.3cm}
  \centering
  \textit{Thank you. We welcome questions --- especially the hard ones.}
\end{frame}

%=================================================
\section*{Acknowledgements}
%=================================================

\begin{frame}{Acknowledgements}
  \small
  \centering

  \vspace{0.5cm}
  {\large \textbf{Thank You}}

  \vspace{0.4cm}
  \textbf{Maverick Team}

  \vspace{0.3cm}
  Junyuan Luo

  \vspace{0.4cm}
  \textit{Sino-Danish College, University of Chinese Academy of Sciences}

  \vspace{0.3cm}
  Digital Innovation Course --- Pitch \& Demo Day

  \vspace{0.3cm}
  10th June, 2026

  \vspace{0.5cm}
  \begin{figure}[htpb]
    \begin{center}
      \includegraphics[width=0.12\linewidth]{pic/CAS_Logo.png}
    \end{center}
  \end{figure}
\end{frame}

\end{document}
